\begin{equation}
    |\psi\rangle_{AB} = |\psi\rangle_A \otimes |\psi\rangle_B
\end{equation}
$|\psi\rangle_{AB}$ : Keadaan produk (\textit{separable}).
Untuk keadaan produk;
\begin{equation}
    \begin{split}
    |\psi \rangle _A &= \frac{1}{\sqrt{2}} (|0\rangle_A + |1\rangle_A) \\
    &= \frac{1}{\sqrt{2}} \begin{pmatrix} 1 \\ 1 \end{pmatrix}
    \end{split}
\end{equation}

\begin{equation}
    \begin{split}
    |\psi \rangle _B &= \frac{1}{\sqrt{2}} (|0\rangle_B + |1\rangle_B) \\
    &= \frac{1}{\sqrt{2}} \begin{pmatrix} 1 \\ 1 \end{pmatrix}
    \end{split}
\end{equation}
maka keadaan produk adalah
\begin{equation}
    \begin{split}
    |\Psi\rangle_{AB} &= |\psi\rangle_A \otimes |\psi\rangle_B \\
    &= \frac{1}{2} (|00\rangle + |01\rangle + |10\rangle + |11\rangle) \\
    &= \frac{1}{2} \begin{pmatrix} 1 \\ 1 \end{pmatrix} \otimes \begin{pmatrix} 1 \\ 1 \end{pmatrix} \\
    &= \frac{1}{2} \begin{pmatrix} 1 \\ 1 \\ 1 \\ 1 \end{pmatrix}
    \end{split}
\end{equation}

sedangkan untuk keadaan terbelit (bell);
\begin{equation}
    |\Psi\rangle_{AB} \ne |\phi\rangle_A \otimes |\phi\rangle_B
\end{equation}
$|\Psi\rangle_{AB}$ : Keadaan terbelit (\textit{entangled}).

di mana
\begin{equation}
    \begin{split}
    |\phi\rangle_A &= \frac{1}{\sqrt{2}} (|0\rangle_A + |1\rangle_A) \\
    &= \frac{1}{\sqrt{2}} \begin{pmatrix} 1 \\ 1 \end{pmatrix}
    \end{split}
\end{equation}
\begin{equation}
    \begin{split}
    |\phi\rangle_B &= \frac{1}{\sqrt{2}} (|0\rangle_B + |1\rangle_B) \\
    &= \frac{1}{\sqrt{2}} \begin{pmatrix} 1 \\ 1 \end{pmatrix}
    \end{split}
\end{equation}
maka
\begin{equation}
    \begin{split}
    |\Phi^+\rangle &= \frac{1}{\sqrt{2}} \left( |00\rangle + |11\rangle \right) \\
    &= \frac{1}{\sqrt{2}} \begin{pmatrix} 1 \\ 0 \\ 0 \\ 1 \end{pmatrix}
    \end{split}
\end{equation}

\begin{equation}
    \rho = \sum_i p_i |\psi_i\rangle \langle \psi_i |
\end{equation}
$\rho$ : Matriks densitas.

\begin{equation}
    \begin{split}
    \rho_{AB} &= |\psi\rangle_{AB} \langle \psi|_{AB} \\
    &= \frac{1}{2} \begin{pmatrix} 1 \\ 1 \end{pmatrix} \frac{1}{2} \begin{pmatrix} 1 & 1 \end{pmatrix} \\
    &= \frac{1}{4} \begin{pmatrix} 1 & 1 \\ 1 & 1 \end{pmatrix}
    \end{split}
\end{equation}

\begin{equation}
    \rho_A = \text{Tr}_B(\rho_{AB})
\end{equation}
$\text{Tr}_B$ : \textit{Partial trace} terhadap subsistem B.

\begin{equation}
    \begin{split}
    \rho_{AB} &= \begin{pmatrix} \rho_{11} & \rho_{12} \\ \rho_{21} & \rho_{22} \end{pmatrix} \\ 
    \end{split}
\end{equation}

\begin{equation}
    \rho_{AB} = \begin{pmatrix} \rho_{00:00} & \rho_{00:01} & \rho_{00:10} & \rho_{00:11} \\ \rho_{01:00} & \rho_{01:01} & \rho_{01:10} & \rho_{01:11} \\ \rho_{10:00} & \rho_{10:01} & \rho_{10:10} & \rho_{10:11} \\ \rho_{11:00} & \rho_{11:01} & \rho_{11:10} & \rho_{11:11} \end{pmatrix}
\end{equation}

\begin{equation}
    \rho_A = \rho_{11} + \rho_{22}
\end{equation}

\begin{equation}
    \rho_A = \begin{pmatrix} \rho_{00:00} + \rho_{10:01} & \rho_{00:01} + \rho_{10:11} \\ \rho_{01:00} + \rho_{11:10} & \rho_{01:01} + \rho_{11:11} \end{pmatrix}
\end{equation}

\begin{equation}
    \begin{split}
    \text{Tr}_B(\rho_{AB}) &= \sum_{i,j} \left(\sum_{k} C_{i,j,k,l} \right) |i \rangle _A \langle j |_A \\
    \end{split}
\end{equation}
untuk keadaan produk $|\Psi\rangle_{AB} = \frac{1}{2}|00\rangle + \frac{1}{2}|01\rangle + \frac{1}{2}|10\rangle + \frac{1}{2}|11\rangle$ 

\begin{equation}
    \rho_{AB} = \frac{1}{4} \begin{pmatrix} 1 & 1 & 1 & 1 \\ 1 & 1 & 1 & 1 \\ 1 & 1 & 1 & 1 \\ 1 & 1 & 1 & 1 \end{pmatrix}
\end{equation}

\begin{equation}
    \begin{split}
    \rho_{A} &= \begin{pmatrix} \frac{1}{4} + \frac{1}{4} & \frac{1}{4} + \frac{1}{4} \\ \frac{1}{4} + \frac{1}{4} & \frac{1}{4} + \frac{1}{4} \end{pmatrix} \\
    &= \begin{pmatrix} \frac{1}{2} & \frac{1}{2} \\ \frac{1}{2} & \frac{1}{2} \end{pmatrix}
    \end{split}
\end{equation}

untuk keadaan bell $|\Phi^+\rangle = \frac{1}{\sqrt{2}} (|00\rangle + |11\rangle)$
\begin{equation}
    \rho_{AB} = \frac{1}{2} \begin{pmatrix} 1 & 0 & 0 & 1 \\ 0 & 0 & 0 & 0 \\ 0 & 0 & 0 & 0 \\ 1 & 0 & 0 & 1 \end{pmatrix}
\end{equation}

\begin{equation}
    \begin{split}
    \rho_{A} &= \begin{pmatrix} \frac{1}{2} + 0 & 0 + 0 \\ 0 + 0 & \frac{1}{2} + 0 \end{pmatrix} \\
    &= \begin{pmatrix} \frac{1}{2} & 0 \\ 0 & \frac{1}{2} \end{pmatrix}
    \end{split}
\end{equation}

ENTROPI SHANNON
\begin{equation}
    H(P) = -\sum_i p_i \log_2 p_i
\end{equation}

\begin{equation}
    S(\rho) = -\text{Tr}(\rho \ln \rho)
\end{equation}
$S(\rho)$ : Entropi Von Neumann.

Dekomposisi Spektral
\begin{equation}
    \rho = U \Lambda U^\dagger
\end{equation}
$\Lambda$ : matriks diagonal real. dengan bentuk $\Lambda = diag(\lambda_1, \lambda_2, ..., \lambda_n)$.

reduksi
\begin{equation}
    \log_2(\rho) = U \log_2(\Lambda) U^\dagger
\end{equation}
\begin{equation}
    \log_2(\Lambda) = diag(\log_2(\lambda_1), \log_2(\lambda_2), ..., \log_2(\lambda_n))
\end{equation}

substitusi
\begin{equation}
    \begin{split}
    S(\rho) &= -\text{Tr}\left[\rho \log_2(\rho)\right] \\
    &= -\text{Tr}\left[ U \Lambda U^\dagger U \log_2(\Lambda) U^\dagger \right] \\
    &= -\text{Tr}\left[ U \Lambda \log_2(\Lambda) U^\dagger \right] \\
    &= -\text{Tr}\left[ \Lambda \log_2(\Lambda) \right] \\
    &= -\text{Tr}\left[ diag(\lambda_1, \lambda_2, ..., \lambda_n) diag(\log_2(\lambda_1), \log_2(\lambda_2), ..., \log_2(\lambda_n)) \right] \\
    &= -\text{Tr}\left[ diag(\lambda_1 \log_2(\lambda_1), \lambda_2 \log_2(\lambda_2), ..., \lambda_n \log_2(\lambda_n)) \right] \\
    &= -\sum_i \lambda_i \log_2(\lambda_i)
    \end{split}
\end{equation}

evaluasi numerik untuk keadaan produk:

\begin{equation}
    \rho_A = \begin{pmatrix} 1/2 & 1/2 \\ 1/2 & 1/2 \end{pmatrix}
\end{equation}

\begin{equation}
    \lambda = \{1, 0\}
\end{equation}

\begin{equation}
    \begin{split}
    S(\rho_A) &= -(1 \log_2(1) + 0 \log_2(0)) \\
    &= 0
    \end{split}
\end{equation}


evaluasi numerik untuk keadaan bell:

\begin{equation}
    \rho_A = \begin{pmatrix} 1/2 & 0 \\ 0 & 1/2 \end{pmatrix}
\end{equation}

\begin{equation}
    \lambda = \{1/2, 1/2\}
\end{equation}

\begin{equation}
    \begin{split}
    S(\rho_A) &= -(1/2 \log_2(1/2) + 1/2 \log_2(1/2)) \\
    &= -\left( -1/2 - 1/2 \right) \\
    &= 1
    \end{split}
\end{equation}
